\chapter{One-Hot Encoding Strategy}\label{Appendix:One_Hot_Encoding}
The one-hot encoding strategy used in this project created $n$ spatial channels based on the number of individual classes passed on to the model. Where traditional one-hot encodiding strategies map the values in a channel to either 1 or 0, referring to the presence or absence of a specific class, the final \textit{Tanh} activation function in the \gls{fluvgan} model maps values to either -1 or 1. Figures~\ref{appendix_fig:one_hot_encoding} and~\ref{appendix_fig:one_hot_encoding_3d} provide an overview of one hot encoding of a vector.

\begin{figure}[htbp!]
    \centering
    \includegraphics[width=0.5\textwidth]{figures/appendix/appendix_B/one_hot_encoding.png}
    \caption{Simpel one-ho encoding represenation by converting the original data array into a multitude of bipolar vectors. One-hot encoding can be applied to any type of categorical data.}
    \label{appendix_fig:one_hot_encoding}
\end{figure}

Using one-hot encoding categorical independent data can be correctly fed to the model, avoiding issues where the model think the numbers/values are correlated. The model may start to correlate the the classes based on values, tending to overestimating the classes that do not occur often.

\begin{figure}[htbp!]
    \centering
    \includegraphics[width=0.99\textwidth]{figures/appendix/appendix_B/one_hot_encoding_3d.png}
    \caption{3-dimensional representation of one-hot encoding used in this project where a 3D reservoir object is subdivided into three separate 3D reservoirs, each difining where a certain facies is or is not present.}
    \label{appendix_fig:one_hot_encoding_3d}
\end{figure}